% ============================================================
%  第五卷 第十章 化学反应（§101–§105 + 习题）
%  底本：高教社中文版第5版（束仁贵/束莼 译，郑伟谋 校）
%  OCR 底稿 full_merged.md 行 10542–10989；对照 原书转换/pages/p-282…p-293.png 核验
%  插图：本章无插图（原书 §101–§105 无图）。
%  §103 节名"反应热"OCR 目录缺行，按正文实况（p271）补排。
%  正文圈码 5 处原样保留（§101 一处、§102 习题1 一处、§104 一处、§105 两处），
%  脚注文本由脚注代理后续转录。
%
%  已修正的 OCR/原书错误（均对照原书页面 PNG 核验）：
%   [§101] 热力学势 Φ（OCR \phi、\mathcal{P} 混用）；ν_i 乱粗体还原为斜体；
%          "2H₂+O₂=2H₂O 即 2H₂+O₂−2H₂O=0"（OCR 0₂、2H₂Λ₀、脱句号）；
%          "平衡条件是由几个(101.2)型的方程式所组成的方程组"一句 OCR 串入
%          "μ_i"字样（p283 无）；∂Φ/∂N_i=μ_i（OCR μ_）；"如果 N₁ 改变 ν₁"
%          （OCR N_\mathfrak{r}）；"化学平衡.一般来讲"（OCR 逗号，p282 为句号）。
%   [§102] 平衡常数下标 K_p（OCR 多处作 K_ρ、K_s）；(102.4) 因子 P^{−Σν_i}
%          （OCR 乱码）；(102.6) 指数 Σc_{pi}ν_i（OCR ε_{pi}）；(102.7)
%          ψ_i(P,T)（OCR ψ_）；习题1 ζ₂ 式尾衍 \quad 删；
%          习题公式 (1)(2)(3) 原书自成编号，自动编号无法产生，按 ch04 (42.5a)
%          体例用 \tag 排。
%   [§102 习题2] (3) 式末"+ 常数"（OCR \frac{\ast}{\ast}***，p286 实为"常数"）；
%          μ₁ 式后 P₁=2α/(1+α)·P——原书排印误作"P=…"（p286），按文意补下标 1；
%          (∂α/∂T)_p（OCR 下标 r）；末式左端 C_p（OCR C_ν，p286 实为 C_p）。
%   [§102 习题3] 答案 c=常数×√P（OCR 乱码 \sqrt{5}/\sqrt{10}×5\sqrt{5}，p286）；
%          反应式 H₂=2H（OCR H₂=2π）。
%   [§103] δ（OCR 8、S、\mathfrak{E} 混用）：δΦ、δn、δN_i、δQ_p、δQ_v、δV；
%          δN_i=−ν_iδn（OCR 乱式）；(103.2) 整式按 p288 重建
%          （OCR \hat{δ}\phi=\frac{T\hat{\sigma}}{F}…）；K_p(T)（OCR K_s）；
%          ∏P_i^{ν_i}>K_p 补 ∏（OCR 脱）；(103.6) 右端 TδnΣν_i（OCR "78n"），
%          前式 δQ_v=δQ_p−δ(PV)（OCR 两处皆作 δQ_ν）；
%          (103.8) 前 δV=(∂/∂P)δΦ（OCR 分母作 δP）；"任意δn."、"完全一致."
%          句末句号补（OCR 脱，p289）。
%   [§104] e⁻（OCR \tt{e^-}）；化学常数 ζ 式补方括号；g=(2L+1)(2S+1)
%          （OCR 倒序）；(104.4) exp(I_n/T)（OCR 分数脱 \frac）。
%   [§105] "的产生"（OCR y^{±}生）；"大得多②"（OCR \Re^{®}）；μ⁺（OCR μ^*）；
%          ε=c√(p²+m²c²)（OCR c\sqrt{p²+m²}c²）；(105.3) 0.183；
%          习题末式 n⁺=n⁻−n₀=−n₀/2+[…]^{1/2}（OCR "n^*=n^-=n_0=…"，
%          p292 实为 n⁺=n⁻−n₀）；"因而"（OCR \overline{\sqrt{\eta}} 乱码）。
%
%  遗留问题：
%   - 脚注仅存圈码（§101、§102 习题1、§104、§105 各见上），脚注文本
%     （p282/p285/p290/p292）由脚注代理转录。
%   - §101 末"条件(101.2)仍旧保持同样依据的形式.这是因为依据相平衡条件"
%     原书如此（p283），照排。
%
%  脚注转录记录（脚注代理，2026-10-12）：
%   - 全章 5 处圈码已转录为 \footnote：p282（§101①）、p285（§102 习题1①）、
%     p290（§104①）、p291（§105①②）。§105 两则实在 p291 页底（原备案 p292 系 ±1 偏差）。
%   - [疑难] §105 脚注②"玻尔半径的立方"原书排印作 (ħ²/me²)³（p291；按物理应为
%     (ħ²/mec²)³），照原书页面转录，未改动，留待校核。
% ============================================================
\chapter{化学反应}

\section{化学平衡条件}

在可反应物质的混合物中进行的化学反应，终将导致建立起一个平衡状态，在其中参与反应的每一种物质的量不再发生变化.这种情况的热平衡称为化学平衡.一般来讲，每一种化学反应都可以在正、逆两个方向进行；在达到平衡以前，两个反应方向中的一个占优势，而在平衡时，这两个相反的反应依适当速率进行，以使得每一种反应物质的粒子总数保持不变.热力学应用于化学反应时，研究的主题只是化学平衡，而不是导致这种平衡的反应过程本身.

重要的是，化学平衡状态不依赖于反应以怎样的方式（以及在怎样的条件下）进行\footnote{尤其是，它与反应是否有催化剂参与无关.}；它只依赖于反应物质的混合物实现平衡时所处的条件.因此在推导化学平衡条件时，可以对反应以怎样的方式进行作任何假设.

首先我们来规定描述化学反应的方法.大家知道，化学反应是以符号方程式的形式来描述的，在普遍形式下可写成（把所有各项都移到方程式一边）：
\begin{equation}
  \sum_i\nu_iA_i=0,
\end{equation}
式中 $A_i$ 是反应物质的化学符号，而系数 $\nu_i$ 是正或负的整数.例如，对于 $2\mathrm{H}_2+\mathrm{O}_2=2\mathrm{H}_2\mathrm{O}$ 即 $2\mathrm{H}_2+\mathrm{O}_2-2\mathrm{H}_2\mathrm{O}=0$ 这个反应来讲，系数 $\nu_{\mathrm{H}_2}=2$，$\nu_{\mathrm{O}_2}=1$，$\nu_{\mathrm{H}_2\mathrm{O}}=-2$.

假设反应在恒温、恒压下进行.在这样的过程中，系统的热力学势趋向于极小值.因此在平衡状态下，热力学势 $\Phi$ 必须具有最小的可能值（在给定的 $P$ 和 $T$ 下）.我们用 $N_1,N_2,\cdots$ 来代表参加反应的各种物质的粒子数.于是 $\Phi$ 为极小值的必要条件可以写成 $\Phi$ 对 $N_i$ 之一如 $N_1$（在给定的 $P$ 和 $T$ 下）的全导数等于零的形式：
\begin{equation*}
  \frac{\partial\Phi}{\partial N_1}
  +\frac{\partial\Phi}{\partial N_2}\frac{\mathrm{d}N_2}{\mathrm{d}N_1}
  +\frac{\partial\Phi}{\partial N_3}\frac{\mathrm{d}N_3}{\mathrm{d}N_1}
  +\cdots=0.
\end{equation*}
各种粒子数 $N_i$ 在反应时的变化，由反应方程式彼此相联系：显然，如果 $N_1$ 改变 $\nu_1$，那么其余每一种粒子数 $N_i$ 各改变 $\nu_i$.换句话说，可以写成 $\mathrm{d}N_i=\dfrac{\nu_i}{\nu_1}\mathrm{d}N_1$，亦即 $\dfrac{\mathrm{d}N_i}{\mathrm{d}N_1}=\dfrac{\nu_i}{\nu_1}$.因此上述等式可以改写成形式
\begin{equation*}
  \sum_i\frac{\partial\Phi}{\partial N_i}\frac{\nu_i}{\nu_1}=0.
\end{equation*}
最后，把 $\dfrac{\partial\Phi}{\partial N_i}=\mu_i$ 代入并消去 $\nu_1$，我们就得到：
\begin{equation}
  \sum_i\nu_i\mu_i=0.
\end{equation}
这就是我们所求的化学平衡条件.因此，要把它写出来，只需在化学反应方程式中把符号 $A_i$ 替换为相应的 $\mu_i$.当混合物中可能有几种不同的反应时，平衡条件是由几个(101.2)型的方程式所组成的方程组.每一个方程式都是用上述方法根据每一种可能反应的方程式所写出的.

应当指出：在反应物质以溶质的形式分配于相互接触的两种不同相中时，条件(101.2)仍旧保持同样依据的形式.这是因为依据相平衡条件，平衡时每一种物质在两相中的化学势彼此相等.

\section{质量作用定律}

我们把在上节中所得到的化学平衡普遍条件应用于在气体混合物中进行的反应，并且假设：气体可以看成是理想气体.

混合物中每一种气体的化学势等于（参看§93）
\begin{equation}
  \mu_i=T\ln P_i+\chi_i(T),
\end{equation}
式中 $P_i$ 是混合物中第 $i$ 种气体的分压强：$P_i=c_iP$.在这里 $P$ 是混合物的总压强，而 $c_i=\dfrac{N_i}{N}$ 是该种气体的浓度，我们把它定义为该种气体的分子数 $N_i$ 对混合物中分子总数 $N=\sum\limits_iN_i$ 之比.

现在很容易写出在气体混合物中反应的化学平衡条件.把(102.1)代入(101.2)，我们就求出：
\begin{equation*}
  \sum_i\nu_i\mu_i=T\sum_i\nu_i\ln P_{0i}+\sum_i\nu_i\chi_i=0,
\end{equation*}
式中 $P_{0i}$ 是各种气体在化学平衡状态下的分压强.引入符号
\begin{equation}
  K_p(T)=\exp\left(-\frac{\sum\nu_i\chi_i}{T}\right),
\end{equation}
我们就得到：
\begin{equation}
  \prod_iP_{0i}^{\nu_i}=K_p(T).
\end{equation}
可以用 $Pc_{0i}$ 来代替 $P_{0i}$，其中 $c_{0i}$ 是气体在化学平衡下的浓度；于是我们得到：
\begin{equation}
  \prod_ic_{0i}^{\nu_i}=P^{-\sum\nu_i}K_p(T)\equiv K_c(P,T).
\end{equation}
在等式(102.3)或(102.4)右边的量只是温度和压强的函数，而与各种反应气体的初始含量无关；这个量通常称为化学平衡常数.而公式(102.3)或(102.4)表示的规律称为质量作用定律.

气体反应平衡常数与压强的关系完全由等式(102.4)右边的因子 $P^{-\sum\nu_i}$ 来确定（如果反应物质的数量由它们的分压强来表示，那么平衡常数就始终与压强无关）.要建立平衡常数对温度的关系，需要对气体的性质作进一步的假设.

例如，如果气体具有恒定的比热，那么把(102.1)式同这种气体的热力学势公式(43.4)比较，就可以证明函数 $\chi_i(T)$ 具有形式
\begin{equation}
  \chi_i(T)=\varepsilon_{0i}-c_{pi}T\ln T-T\zeta_i,
\end{equation}
式中 $c_{pi}$ 是气体的比热，$\zeta_i$ 是气体的化学常数.把这个式子代入(102.2)，我们就得到平衡常数的如下公式：
\begin{equation}
  K_p(T)=\mathrm{e}^{\sum\nu_i\zeta_i}T^{\sum c_{pi}\nu_i}
  \mathrm{e}^{-\sum\nu_i\varepsilon_{0i}/T}.
\end{equation}
它对温度的依赖关系主要是指数规律.

质量作用定律对于溶质之间的反应也成立，只要溶液可以认为是稀溶液.事实上，每一种溶质的化学势具有形式
\begin{equation}
  \mu_i=T\ln c_i+\psi_i(P,T).
\end{equation}
浓度 $c_i$ 在这里被定义为该种溶质的粒子数与溶剂粒子数之比 $\left(c_i=\dfrac{n_i}{N}\right)$.把(102.7)代入平衡条件(101.2)，可以用同样的方法求出：
\begin{equation}
  \prod_ic_{0i}^{\nu_i}=K(P,T),
\end{equation}
式中平衡常数
\begin{equation}
  K(P,T)=\exp\left(-\frac{\sum\nu_i\psi_i}{T}\right).
\end{equation}
与气体反应的情形不同，在这里平衡常数对压强的依赖关系仍旧是未定的.

如果除了气体和溶质以外，还有处于纯凝聚相如纯固体（即没有同其它物质混合）的某种物质参加反应，那么平衡条件也会导致质量作用定律.但是这时因为纯相的化学势只与压强和温度有关，所以在这个定律的方程式左边不包含有关纯相的因子，即只需要写出气体（或溶质）浓度的乘积，就好像固体根本不存在一样.固体的存在只影响平衡常数对压强和温度的依赖关系.

如果参加反应的只有气体和固体，则由于气体的压强比较小，可认为固体的化学势与压强无关，因而平衡常数与压强的关系仍旧像在(102.4)中一样.但是这时指数中的和式 $\sum\nu_i$ 必须只表示反应方程式中气态物质的系数之和.

最后，如果稀溶液中除了溶质以外，溶剂也参加反应，那么这时质量作用定律也成立.事实上，当把溶剂的化学势代入化学平衡条件时，可把溶剂化学势中包含浓度的小项忽略不计，于是溶剂的化学势化为一个只与温度和压强有关的量.因此我们再次得到质量作用定律的方程式，并且在这个方程式的左边也只包含被溶反应物质的浓度，而不包含溶剂项.

\begin{xiti}

\noindent{\heiti 习题1}\quad 设气体的分子由相同双原子构成，并且在基态没有自旋和轨道角动量，试求双原子气体在高温下的离解度.

解：所要讨论的是 $\mathrm{A}_2=2\mathrm{A}$ 型的反应.我们将（在本题和下列各题中）用指标 1 和 2 分别标记混合物中属于原子（$\mathrm{A}$）和分子（$\mathrm{A}_2$）组分的量.离解度定义为已离解分子数 $\dfrac{N_1}{2}$ 与（不可离解气体该有的）分子总数 $N_0=N_2+\dfrac{N_1}{2}$ 之比 $\alpha=\dfrac{N_1}{2N_0}$.根据质量作用定律(102.3)有
\begin{equation*}
  \frac{P_2}{P_1^2}=\frac{N_2(N_1+N_2)}{PN_1^2}
  =\frac{1-\alpha^2}{4\alpha^2P}=K_p(T).\tag{1}
\end{equation*}
由此可得
\begin{equation*}
  \alpha=\frac{1}{\sqrt{1+4PK_p(T)}}.
\end{equation*}
化学平衡常数 $K_p$ 可由将以下各量代入(102.6)式求得：比热 $c_{p_1}=\dfrac52$，$c_{p_2}=\dfrac92$ 以及化学常数
\begin{equation*}
  \zeta_1=\ln\left[g_1\left(\frac{m}{2\pi\hbar^2}\right)^{3/2}\right],
  \qquad
  \zeta_2=\ln\left[\frac{I}{\hbar^6\omega}
  \left(\frac{m}{\pi}\right)^{3/2}\right]
\end{equation*}
（参看(45.4)，(46.4)，(49.8)），其中 $m$ 是原子 $\mathrm{A}$ 的质量，$g_1$ 是原子 $\mathrm{A}$ 基态的统计权重（在足够高的温度下 $g_1=(2S+1)(2L+1)$，这里 $S,L$ 是原子的自旋和轨道角动量）\footnote{参看§104的脚注.}.结果为
\begin{equation*}
  K_p(T)=\frac{8I\pi^{3/2}}{g_1^2\omega m^{3/2}T^{1/2}}\,
  \mathrm{e}^{\varepsilon_0/T}.\tag{2}
\end{equation*}
式中 $\varepsilon_0=2\varepsilon_{01}-\varepsilon_{02}$ 是分子的离解能.

\noindent{\heiti 习题2}\quad 试求上题中可离解双原子气体的热容，分子的离解能为 $\varepsilon_0=2\varepsilon_{01}-\varepsilon_{02}$.

解：气体的熵由以下和式计算：
\begin{equation*}
\begin{aligned}
  S&=N_1\left(c_{p1}+\frac{\varepsilon_{01}-\mu_1}{T}\right)
  +N_2\left(c_{p2}+\frac{\varepsilon_{02}-\mu_2}{T}\right)=\\
  &=N_1\left(c_{p1}+\frac{\varepsilon_{01}}{T}\right)
  +N_2\left(c_{p2}+\frac{\varepsilon_{02}}{T}\right)
  -(N_1+2N_2)\frac{\mu_1}{T}
\end{aligned}
\end{equation*}
（每一个组分的熵，根据(43.6)和(43.3)用其化学势表示，之后用了平衡方程 $\mu_2=2\mu_1$）.用 $N_0$ 和 $\alpha$ 表示 $N_1$ 和 $N_2$，把化学势写成
\begin{equation*}
  \mu_1=\varepsilon_{01}+T\ln P_1-c_{p1}T\ln T-\zeta_1T,
  \qquad P_1=\frac{2\alpha}{1+\alpha}P
\end{equation*}
并且把 $c_{p1}$ 和 $c_{p2}$ 的值代入，我们得到
\begin{equation*}
  S=N_0\left(-\frac{1-\alpha}{T}\varepsilon_0+5\ln T+\frac{\alpha}{2}
  -2\ln\frac{\alpha}{1+\alpha}+\text{常数}\right),\tag{3}
\end{equation*}
这里也用到离解能 $\varepsilon_0$，而这里的常数是温度无关项，对所求的热容 $C_p=T\left(\dfrac{\partial S}{\partial T}\right)_p$ 没有影响.从(1)算出导数
\begin{equation*}
  \left(\frac{\partial\alpha}{\partial T}\right)_p
  =-\frac{(1-\alpha^2)\alpha}{2}\frac{\mathrm{d}\ln K_p}{\mathrm{d}T}
  =\frac{(1-\alpha^2)\alpha}{2T}
  \left(\frac{\varepsilon_0}{T}+\frac12\right)
\end{equation*}
（$K_p$ 取(2)式）.现在对(3)的熵求微商，最后我们得到
\begin{equation*}
  C_p=\frac{N_0}{2}\left[9+\alpha+\alpha(1-\alpha^2)
  \left(\frac{\varepsilon_0}{T}+\frac12\right)^2\right].
\end{equation*}

\noindent{\heiti 习题3}\quad 试确定氢以 $\mathrm{H}$ 原子形式溶解在金属中的浓度对金属上的 $\mathrm{H}_2$ 气压的依赖关系.

解：把过程考虑为化学反应 $\mathrm{H}_2=2\mathrm{H}$，平衡条件可写成形式 $\mu_{\mathrm{H}_2}=2\mu_{\mathrm{H}}$；我们把 $\mu_{\mathrm{H}_2}$ 写成理想气体的化学势：
\begin{equation*}
  \mu_{\mathrm{H}_2}=T\ln P+\chi(T),
\end{equation*}
把 $\mu_{\mathrm{H}}$ 写成溶质在溶液中的化学势：
\begin{equation*}
  \mu_{\mathrm{H}}=T\ln c+\psi.
\end{equation*}
注意到 $\psi$ 几乎与压强无关（参看§90），我们就求出：
\begin{equation*}
  c=\text{常数}\times\sqrt{P}.
\end{equation*}

\end{xiti}

\section{反应热}

化学反应都伴有热量的吸收或放出.前一种情形称为吸热反应，后一种情形称为放热反应.显然，如果某种反应是放热的，那么它的逆反应就是吸热的，反之亦然.

反应的热效应与发生反应的条件有关.因此，在讨论热效应时必须区分譬如说反应是在恒定的体积还是压强下进行的.（但是这种差别通常不大.）

像计算溶解热的情形（§91）一样，我们首先来确定由于化学反应可能获得的最大功.

我们把由一个化学反应方程式所确定的一组分子之间的反应称为“元反应”，并且来计算反应物质混合物在元反应发生了很少的 $\delta n$ 次后热力学势的改变；这里我们假设反应在恒温恒压下进行.于是我们得到：
\begin{equation*}
  \delta\Phi=\sum_i\frac{\partial\Phi}{\partial N_i}\delta N_i
  =\sum_i\mu_i\delta N_i.
\end{equation*}
在 $\delta n$ 次的元反应后第 $i$ 种物质分子数的改变显然等于
\begin{equation*}
  \delta N_i=-\nu_i\delta n.
\end{equation*}
因此，
\begin{equation}
  \delta\Phi=-\delta n\sum_i\nu_i\mu_i.
\end{equation}
由此可见，在平衡状态下 $\dfrac{\delta\Phi}{\delta n}$ 变成0，这原是理所当然的.

式(103.1)是为了进行 $\delta n$ 次的元反应所必须耗费的最小功的普遍表达式.它也是同样次数元反应在逆方向发生时可能获得的最大功.

首先我们假设反应发生在气体之间.利用 $\mu_i$ 的表达式(102.1)，我们就求出：
\begin{equation*}
  \delta\Phi=-\delta n\left(T\sum_i\nu_i\ln P_i
  +\sum_i\nu_i\chi_i\right),
\end{equation*}
引入平衡常数后，得到：
\begin{equation}
  \delta\Phi=T\delta n\left[-\sum_i\nu_i\ln P_i+\ln K_p(T)\right]
  =T\delta n\left[-\sum_i\nu_i\ln c_i+\ln K_c(P,T)\right].
\end{equation}
对于在溶液中进行的反应，我们可借助于 $\mu_i$ 的表达式(102.7)类似地求出：
\begin{equation*}
  \delta\Phi=-\delta n\left(T\sum_i\nu_i\ln c_i
  +\sum_i\nu_i\psi_i\right),
\end{equation*}
引入平衡常数 $K(P,T)$ 后，得到：
\begin{equation}
  \delta\Phi=T\delta n\left[-\sum_i\nu_i\ln c_i
  +\ln K(P,T)\right].
\end{equation}
值 $\delta\Phi$ 的正负决定反应进行的方向：因为 $\Phi$ 趋向于极小值，所以当 $\delta\Phi<0$ 时反应沿正方向（即化学反应方程式中“从左向右”的方向）进行，如果 $\delta\Phi>0$ 则意味着反应在所给定的混合物中实际上沿逆方向进行.但是，应当注意，反应方向也可以从质量作用定律直接看出：我们对所给定的混合物写下乘积 $\prod\limits_iP_i^{\nu_i}$，并同该反应的平衡常数值进行比较；例如，如果是 $\prod\limits_iP_i^{\nu_i}>K_p$，那么这就意味着反应将沿正方向进行，以使得（反应方程式中 $\nu_i$ 为正的）初始物质的分压强减小，而（$\nu_i<0$ 的）反应产物的分压强增大.

现在我们也可以对 $\delta n$ 次的元反应确定所吸收（或所放出，视符号而定）的热量.根据公式(91.4)，对于在恒温恒压下所进行的反应来讲，这个热量 $\delta Q_p$ 等于
\begin{equation*}
  \delta Q_p=-T^2\left(\frac{\partial}{\partial T}
  \frac{\delta\Phi}{T}\right)_p.
\end{equation*}
对于气体之间的反应，把(103.2)代入，我们就得到：
\begin{equation}
  \delta Q_p=-T^2\delta n\frac{\partial\ln K_p(T)}{\partial T}.
\end{equation}
类似地，对于溶液，
\begin{equation}
  \delta Q_p=-T^2\delta n\frac{\partial\ln K(P,T)}{\partial T}.
\end{equation}
必须注意：$\delta Q_p$ 只是正比于 $\delta n$ 而与任何给定时刻的浓度值无关；因此这些公式也适用于不很小的任意 $\delta n$.

如果 $Q_p>0$，即反应是吸热的，则 $\dfrac{\partial\ln K}{\partial T}<0$，即化学平衡常数随着温度的增加而下降.相反地，对于放热反应（$Q_p<0$），平衡常数随温度而增长.从另一方面来讲，平衡常数的增长意味着化学平衡朝着重新形成初始物质的方向移动（反应“从右向左”进行），以使乘积 $\prod\limits_ic_{0i}^{\nu_i}$ 增加.相反地，化学平衡常数的减小意味着化学平衡朝着形成反应产物的方向移动.换句话说，可以表述法则如下：加热使平衡朝着过程伴有吸热的方向移动，而冷却使平衡朝着过程伴有放热的方向移动.这个法则与勒夏特列原理完全一致.

对于气体之间的反应，在恒定体积下（同时在恒温下）进行反应时的热效应也很重要.热量 $\delta Q_v$ 同热量 $\delta Q_p$ 之间的关系很简单.事实上，在恒定体积过程中所吸收的热量等于系统能量的改变，而 $\delta Q_p$ 则等于焓的改变.因为 $E=W-PV$，显然
\begin{equation*}
  \delta Q_v=\delta Q_p-\delta(PV),
\end{equation*}
或者再代入 $PV=T\sum N_i$ 和 $\delta N_i=-\nu_i\delta n$，得
\begin{equation}
  \delta Q_v=\delta Q_p+T\delta n\sum_i\nu_i.
\end{equation}

最后，我们再来确定恒压（和恒温）下进行反应所导致的反应物质混合物的体积改变.对于气体，这个问题很容易：
\begin{equation}
  \delta V=\frac TP\delta N=-\frac TP\delta n\sum_i\nu_i.
\end{equation}
特别地，如果反应不改变粒子总数（$\sum\limits_i\nu_i=0$），那么反应进行时体积也不发生变化.

对于稀溶液中的反应，我们可以利用公式 $\delta V=\dfrac{\partial}{\partial P}\delta\Phi$，并且把(103.3)式代入，就得到：
\begin{equation}
  \delta V=T\delta n\frac{\partial\ln K(P,T)}{\partial P}.
\end{equation}
因此，反应时体积的变化来自化学平衡常数对压强的依赖关系.类似于上述关于温度依赖关系的讨论，很容易得出结论：增加压强将促进伴有体积减小的反应（即把平衡朝这个反应方向移动），而减小压强将促进伴有体积增加的反应——这仍旧与勒夏特列原理完全一致.

\section{电离平衡}

在足够高的温度下，气体粒子的碰撞可伴有电离.这种热电离的存在导致一种热平衡的建立，那里气体粒子以总数的一定比例处于不同的电离阶段.我们来考虑单原子气体的热电离；这种情形最为重要，因为化合物通常在出现热电离之前就已经完全离解了.

从热力学的观点来看，电离平衡乃是化学平衡的特殊情形，它相当于同时发生的可写成如下形式的多个“电离反应”：
\begin{equation}
  \mathrm{A}_0=\mathrm{A}_1+\mathrm{e}^-,\qquad
  \mathrm{A}_1=\mathrm{A}_2+\mathrm{e}^-,\quad\cdots,
\end{equation}
式中符号 $\mathrm{A}_0$ 表示中性原子，$\mathrm{A}_1,\mathrm{A}_2,\cdots$ 表示一次离解、二次离解……的原子，$\mathrm{e}^-$ 表示电子.把质量作用定律应用于这些反应，就得出一组方程组
\begin{equation}
  \frac{c_{n-1}}{c_nc}=PK_p^{(n)}(T)\qquad(n=1,2,\cdots),
\end{equation}
式中 $c_0$ 是中性原子的浓度，$c_1,c_2,\cdots$ 是各次离子的浓度，$c$ 是电子的浓度（这些浓度中的每一个都定义为该种粒子数目与包括电子在内的粒子总数之比）.这些方程式还须加上表示气体作为整体有电中性的方程式
\begin{equation}
  c=c_1+2c_2+3c_3+\cdots.
\end{equation}
方程组(104.2—104.3)确定在电离平衡下各种离子的浓度.

平衡常数 $K_p^{(n)}$ 很容易计算出来.所有参加反应的气体（中性原子气体、离子气体、电子气）都是“单原子”的，因而都具有 $c_p=\dfrac52$ 的恒定比热，而它们的化学常数等于
\begin{equation*}
  \zeta=\ln\left[g\left(\frac{m}{2\pi\hbar^2}\right)^{3/2}\right],
\end{equation*}
式中 $m$ 是该种气体粒子的质量，$g$ 是该种粒子基态的统计权重；对于电子来讲，$g=2$，对于原子和离子，$g=(2L+1)(2S+1)$（$L,S$ 各为原子或离子的轨道角动量和自旋）\footnote{可以认为，即使在充分电离的气体中，所有的原子和离子也都处于基态中；其原因在于：只要原子（或离子）的基态具有精细结构，那么我们总是可以假设 $T$ 比这种结构的能量间距大得多.}.把这些值代入公式(102.6)，我们就得到所求的平衡常数的表达式：
\begin{equation}
  K_p^{(n)}(T)=\frac{g_{n-1}}{2g_n}
  \left(\frac{2\pi}{m}\right)^{3/2}
  \frac{\hbar^3}{T^{5/2}}\exp\frac{I_n}{T}
\end{equation}
（M.Saha，1921），式中 $m$ 是电子的质量，而 $I_n=\varepsilon_{0,n}-\varepsilon_{0,n-1}$ 是原子的第 $n$ 次电离能（第 $n$ 次电离势）.

平衡常数 $K_c^{(n)}=PK_p^{(n)}$ 随着温度的增加而减小，当它达到1的量级时，气体的第 $n$ 次电离的电离度也为1的量级.特别应当注意的是：虽然平衡常数的温度依赖关系是指数型的，但是这个电离阶段并不是在 $T\sim I_n$ 时才发生，而是在低得多的温度下就已经发生了.其原因在于指数因子 $\exp\left(\dfrac{I_n}{T}\right)$ 的系数很小；事实上，
\begin{equation*}
  \frac PT\left(\frac{\hbar^2}{mT}\right)^{3/2}
  =\frac NV\left(\frac{\hbar^2}{mT}\right)^{3/2},
\end{equation*}
这个量一般来讲很小，当 $T\sim I$ 时，它的数量级为原子体积与一个原子在气体中所占体积即 $\dfrac VN$ 之比.

因此，气体在比电离能小得多的温度下就已经充分电离了.但是同时气体中受激原子的数目仍旧非常少，因为一般来讲原子的激发能与电离能同数量级.至于当 $T$ 同电离能可比时，这时气体已经差不多完全电离.当温度数量级达到原子中最后一个电子的电离能时，可以认为气体只由单个的电子和裸核所构成.

第一个电子的电离能 $I_1$ 通常都比其次的各能值 $I_n$ 小得多；因此存在一个温度范围，可以认为那里的气体除了中性原子以外只有带单电荷的离子.我们用电离原子数与原子总数之比 $\alpha$ 作为气体的电离度，于是有：
\begin{equation*}
  c=c_1=\frac{\alpha}{1+\alpha},\qquad
  c_0=\frac{1-\alpha}{1+\alpha},
\end{equation*}
而方程(104.2)给出：
\begin{equation*}
  \frac{1-\alpha^2}{\alpha^2}=PK_p^{(1)},
\end{equation*}
由此
\begin{equation}
  \alpha=\frac{1}{\sqrt{1+PK_p^{(1)}}},
\end{equation}
由这个式子就完全确定了电离度与压强和温度（在所考虑的温度范围内）的关系.

\section{涉及粒子对产生的平衡}

在与电子的静止能量 $mc^2$ 可比的极高温度下\footnote{能量 $mc^{2}=0.51\times10^{6}\,\mathrm{eV}$，因此温度 $\dfrac{mc^{2}}{k}=6\times10^{9}\,\mathrm{K}$.}，物质内粒子的碰撞伴随着电子对（电子和正电子）的产生；由于这个缘故，粒子数不再是一个给定的量，而决定于热平衡条件.

电子对的产生（及其湮没）从热力学的观点而言，可以看成“化学反应”
\begin{equation*}
  \mathrm{e}^++\mathrm{e}^-=\gamma,
\end{equation*}
式中符号 $\mathrm{e}^+$ 和 $\mathrm{e}^-$ 代表正电子和电子，而符号 $\gamma$ 代表一个或几个光子.光子气体的化学势等于0（§63）.因此，产生偶对的平衡条件如下式：
\begin{equation}
  \mu^-+\mu^+=0,
\end{equation}
式中 $\mu^-$ 和 $\mu^+$ 各为电子气和正电子气的化学势.应当强调指出：在这里 $\mu$ 指化学势的相对论表达式，它包含粒子的静止能量（参看§27），后者为参与粒子对产生过程的关键角色.

在 $T\sim mc^2$ 的温度下，（在单位体积内）产生的电子对的数目就已经比原子的电子密度大得多\footnote{从公式(105.3)可以看出，当 $T\sim mc^{2}$ 时，生成的电子对的体积为 $\sim(\hbar/mc)^{3}$. 比起原子体积即玻尔半径的立方 $(\hbar^{2}/me^{2})^{3}$，这个体积非常小.}.因此可以足够精确地认为：电子数等于正电子数.于是 $\mu^-=\mu^+$，因而条件(105.1)给出：
\begin{equation*}
  \mu^-=\mu^+=0,
\end{equation*}
亦即在平衡状态下电子和正电子的化学势必须都等于0.

电子和正电子服从费米统计；因此把 $\mu=0$ 的费米分布(56.3)进行积分，就得到它们的数目：
\begin{equation}
  N^+=N^-=\frac{V}{\pi^2\hbar^3}
  \int_0^{\infty}\frac{p^2\,\mathrm{d}p}{\mathrm{e}^{\varepsilon/T}+1},
\end{equation}
式中 $\varepsilon$ 由相对论表达式 $\varepsilon=c\sqrt{p^2+m^2c^2}$ 来确定.

当 $T\ll mc^2$ 时，这个数目正比于指数函数 $\exp\left(-\dfrac{mc^2}{T}\right)$，因而很小.在相反的情形下，$T\gg mc^2$.我们可以令 $\varepsilon=cp$，因而公式(105.2)给出：
\begin{equation*}
  N^+=N^-=\frac{V}{\pi^2}\left(\frac{T}{\hbar c}\right)^3
  \int_0^{\infty}\frac{x^2\,\mathrm{d}x}{\mathrm{e}^x+1}.
\end{equation*}
这里的积分可以用 $\zeta$ 函数（见§58第二个脚注）来表示，因而我们得到：
\begin{equation}
  N^+=N^-=\frac{3\zeta(3)}{2\pi^2}\left(\frac{T}{\hbar c}\right)^3
  =0.183\left(\frac{T}{\hbar c}\right)^3V.
\end{equation}
用同样的方法可以求出正电子气和电子气的能量：
\begin{equation*}
  E^+=E^-=\frac{VT}{\pi^2}\left(\frac{T}{\hbar c}\right)^3
  \int_0^{\infty}\frac{x^3\,\mathrm{d}x}{\mathrm{e}^x+1}
  =\frac{7\pi^2T^4}{120(\hbar c)^3}V.
\end{equation*}
这个量是在同样体积中黑体辐射能量的 $7/8$.

\begin{xiti}

\noindent{\heiti 习题}\quad 试求出当 $T\ll mc^2$ 时电子和正电子的平衡密度.

解：利用化学势的表达式(46.1a)，并把它加上 $mc^2$，我们就得到：
\begin{equation*}
  n^+n^-=4\left(\frac{mT}{2\pi\hbar^2}\right)^3
  \exp\left(-\frac{2mc^2}{T}\right),
\end{equation*}
式中 $n^-=N^-/V$ 和 $n^+=N^+/V$ 各为电子和正电子的密度.如果 $n_0$ 为电子的初始密度（即没有电子对的产生时的密度），则 $n^-=n^++n_0$，因而我们得到：
\begin{equation*}
  n^+=n^--n_0=-\frac{n_0}{2}
  +\bigg[\frac{n_0^2}{4}
  +\frac{1}{2\pi^3}\left(\frac{mc}{\hbar}\right)^6
  \left(\frac{T}{mc^2}\right)^3
  \exp\left(-\frac{2mc^2}{T}\right)\bigg]^{1/2}.
\end{equation*}

\end{xiti}
